\documentclass[11pt]{article}
\usepackage[T1]{fontenc}
\usepackage[utf8]{inputenc}
\usepackage{amsmath,amssymb,amsthm,booktabs}
\usepackage[margin=1in]{geometry}
\usepackage{hyperref,xurl}
\newtheorem{theorem}{Theorem}
\newtheorem{corollary}[theorem]{Corollary}
\newcommand{\ip}[2]{\langle #1,#2\rangle}
\title{Five-dimensional attachment and incidence constraints\linebreak for right-angled hyperbolic polytopes}
\author{Coxeter Polytopes Frontier Project}
\date{18 October 2026}
\begin{document}
\maketitle
\begin{abstract}
Every finite-volume entirely ideal right-angled hyperbolic four-polytope
occurs as an actual facet of a finite-faced infinite-volume right-angled
chamber in hyperbolic five-space. All its original cusps can simultaneously
have bounded rectangular four-box links. An exact positive-carrier example
has 457 genuine facets and 258 such cusps. This removes a possible
one-layer obstruction, without solving finite-volume completion.
Separately, the internally proved carrier and second-volume equality
theorems imply that an entirely ideal right-angled five-polytope needs at
least 56 facets. Published Nikulin--Khovanskii inequalities then restrict
the fraction of entirely ideal five-faces in dimensions nine to twelve;
in dimensions eleven and twelve it is strictly less than one tenth.
These are necessary constraints, not new dimension exclusions.
\end{abstract}

\section{A universal attachment with genuine facets}
Use the Lorentz metric of signature $(4,1)$ on the old vector space.
Let $Q\subset\mathbb H^4$ be a finite-volume entirely ideal right-angled
polytope. Choose a strict interior vector $p=(1,0,0,0,0)$ and normalize
its ideal vertices as $v=(1,x)$. Write $m_i$ for its outward unit facet
normals, so $m_{i0}>0$. Extend the vector space by an orthogonal unit
spacelike vector $n$.

\begin{theorem}
For sufficiently large $L>1$, the halfspaces with outward normals
\[
 n,\quad m_i,\quad h_v=Lv-n
\]
define a full-dimensional right-angled chamber in $\mathbb H^5$ with
exactly $1+f_3(Q)+f_0(Q)$ genuine facets. Its intersection with $n^\perp$
is exactly $Q$, and every original ideal vertex has a bounded rectangular
four-box cusp link. The chamber has infinite volume.
\end{theorem}
\begin{proof}
Put
\[
 a=\min_{\ip{v}{m_i}<0}(-\ip{v}{m_i})>0,
 \qquad b=\min_{v\ne w}(-\ip{v}{w})>0.
\]
Choose $L>1$ with $La>1$ and $L^2b>2$. Then
\[
 h_v^2=1,\quad \ip{n}{h_v}=-1,\quad
 \ip{h_v}{m_i}\in\{0\}\cup(-\infty,-1),\quad
 \ip{h_v}{h_w}=1+L^2\ip{v}{w}<-1\quad(v\ne w).
\]
All wall pairs are orthogonal, asymptotic or ultraparallel, and
$p-n/2$ is a strict future timelike interior vector.

On $n^\perp$, the old inequalities define $Q$ and every new inequality is
strict at each finite point of $Q$: future timelike and future null
vectors have negative Lorentz product. Thus $Q$ is the actual $n$ facet,
with relative-interior witness $p$.

For an old facet choose a finite relative-interior vector $q$ in $Q$.
For sufficiently small $\epsilon>0$, $q-\epsilon n$ is timelike, lies on
that old wall, and has strictly negative product with all other walls.
For a new wall use
\[
 x_v=v+\frac{p-Ln}{L^2}.
\]
It is future timelike with $x_v^2=-L^{-2}-L^{-4}$, and
\[
 \begin{aligned}
 \ip{x_v}{h_v}&=0,& \ip{x_v}{n}&=-L^{-1},\\
 \ip{x_v}{m_i}&=\ip{v}{m_i}-m_{i0}/L^2<0,&
 \ip{x_v}{h_w}&=L\ip{v}{w}<0\quad(w\ne v).
 \end{aligned}
\]
Every listed wall is consequently genuine.

At each old ideal vertex $v$, precisely eight walls meet: $n,h_v$ and
the six old incident walls. The old walls form three orthogonal opposite
pairs whose sums are positive multiples of $v$; the new pair has sum
$n+h_v=Lv$. The eight normals have rank five. The horospherical
intersection is therefore a bounded rectangular four-box, while all
remaining inequalities are strict near this cusp.

Finally, $u=p-n$ is future null and
\[
 \ip{u}{n}=-1,\qquad \ip{u}{m_i}=-m_{i0}<0,
 \qquad \ip{u}{h_v}=1-L<0.
\]
An open ideal spherical patch around its direction and an inward collar
satisfy all inequalities. The Klein volume density diverges along such
a collar, so the chamber has infinite volume.
\end{proof}

The new opposite normals have the necessary form, rather than being a
choice of unrelated local diagrams. The six old cusp normals span a
rank-four subspace with null radical $\mathbb Rv$; its orthogonal
complement in the old Lorentz space is $\mathbb Rv$. A unit normal
opposite $n$ and orthogonal to all six must therefore be
$\lambda_vv-n$. Its inequality is strict on the interior of $Q$ only
when $\lambda_v>0$. The theorem chooses these parameters simultaneously.

For the frozen sixteen-cell positive-carrier example,
\[
 f_0(Q)=258,\quad f_3(Q)=198,\quad \kappa(Q)=6,
 \quad a=1/25,\quad b=1/1150.
\]
The rational choice $L=48$ works. The exact certificate checks all
104,196 wall pairs, all 457 strict facet-interior witnesses, all 258
complete cusp links, the strict center and the open ideal escape.
No final cone census is needed for these claims. The coordinates are
rational in standard Lorentz space; no finite-covolume arithmetic-group
claim is made. Additional walls may cut the escape, so finite-volume
completion remains open.

\section{A necessary 56-facet bound in the ideal five-dimensional case}
For an ideal right-angled four-polytope define
\[
 \Delta(Q)=f_0(Q)-f_3(Q),\qquad
 \kappa(Q)=3f_0(Q)-4f_3(Q)+24.
\]
The separate internal carrier theorem proves the facet identity
$\sum_Q\kappa(Q)=48$ for an ideal right-angled five-polytope. The separate
internal second-volume and equality classification implies that
$\kappa(Q)>0$ forces $\Delta(Q)\ge5$. These geometric theorems are
dependencies, not conclusions of the rational calculation below.

\begin{theorem}
Every finite-volume entirely ideal right-angled hyperbolic five-polytope
has at least 56 facets.
\end{theorem}
\begin{proof}
Choose a positive-$\kappa$ four-facet $Q$. Kolpakov's published facet
minimum and equality theorem \cite{Kolpakov} gives $f_3(Q)\ge24$, with
equality only for the regular 24-cell. That cell has $\kappa=0$, so
$f_3(Q)\ge25$ and $f_0(Q)=f_3(Q)+\Delta(Q)\ge30$.

There are $f_3(Q)$ distinct ridge-neighbor facets of the ambient
five-polytope. At each of the $f_0(Q)$ cusps of $Q$, the rectangular
four-box link supplies a unique ambient facet opposite $Q$.
Its hyperplane is asymptotic to the hyperplane of $Q$ at precisely that
ideal point. It cannot recur at a different cusp, nor equal an orthogonal
ridge neighbor. Counting $Q$ itself gives
\[
 f_4(P)\ge1+f_3(Q)+f_0(Q)\ge56.
\]
\end{proof}

\section{Published inputs and mixed-dimensional consequences}
A finite-volume right-angled polytope is simple at edges; all its faces
are finite-volume and right-angled. We use the published
Nikulin--Khovanskii strict average inequality
\[
 U(n,k,l)=\binom{n-l}{n-k}
 \frac{\binom{\lfloor n/2\rfloor}{l}+\binom{\lceil n/2\rceil}{l}}
 {\binom{\lfloor n/2\rfloor}{k}+\binom{\lceil n/2\rceil}{k}},
 \qquad
 \frac1{f_k}\sum_{E\in\Omega_k}f_l(E)<U(n,k,l),
\]
for $0\le l<k\le\lceil n/2\rceil$. We use only $k=5$, $l=1,4$,
$n=9,10,11,12$. Alexandrov \cite{Alexandrov} and Egorov
\cite{Egorov} state the formula and hypotheses explicitly.
We also use Dufour's published noncompact bound
$f_4(P^5)\ge1+\nu_4=16$, and the published absence of compact
right-angled five-polytopes. The latter removes any compact exception.
The original Dufour full article was not accessible for independent
proof review; its statements are reported explicitly in the cited
primary preprints. No equality classification from that article is assumed.

Let $F=f_5(P)$ and let $I$ count the entirely ideal five-faces. Every
five-face has at least 16 facets, and the ideal ones have at least 56.
Thus
\[
 16F+40I\le\sum_{E\in\Omega_5(P)}f_4(E)<U(n,5,4)F.
\]
Writing $c=1-(U(n,5,4)-16)/40$, we obtain $F-I>cF$.
Every nonideal five-face contains a finite vertex, and each finite
vertex belongs to $\binom n5$ five-faces. Hence
\[
 \binom n5v_f\ge F-I>cF.
\]
An edge is contained in $\binom{n-1}{4}$ five-faces and
$2f_1=nv_f+2^{n-1}v_i$. Here $2^{n-1}$ is the number of vertices of
the ideal rectangular-box link, not its number of facets. Applying the
edge average inequality gives
\[
 F>a v_f+b v_i,\qquad
 a=\frac{n\binom{n-1}{4}}{2U(n,5,1)},\quad
 b=\frac{2^{n-2}\binom{n-1}{4}}{U(n,5,1)}.
\]
Since $\binom n5-ca>0$, elimination proves
\[
 v_f>\frac{cb}{\binom n5-ca}v_i.
\]
The exact values are
\begin{center}
\begin{tabular}{ccccc}\toprule
$n$&$U(n,5,4)$&$I/F<$&$F>$&$v_f/v_i>$\\\midrule
9&30&$7/20$&$v_f/2+128v_i/9$&$3328/45243$\\
10&30&$7/20$&$v_f+256v_i/5$&$3328/25135$\\
11&20&$1/10$&$7v_f/2+3584v_i/11$&$3072/4807$\\
12&20&$1/10$&$6v_f+1024v_i$&$512/437$\\\bottomrule
\end{tabular}
\end{center}
The ratio notation presumes $v_i>0$; the displayed linear inequality
also applies when $v_i=0$. In dimension twelve the necessary relation is
$437v_f>512v_i$. A boundary with at least one tenth of its five-faces
entirely ideal is forbidden in dimensions eleven and twelve. These
conclusions neither exclude the dimensions nor assert optimality among
all published vertex bounds.

\section{Proof and computation status}
The universal attachment has a full independent internal mathematical
review of its proof, producer and standalone checker. Exact generation,
fresh replay and the standalone checker passed. The density proof has
two independent full internal mathematical reviews, with precise scopes
recorded in its companion report. Its script checks rational
consequences of explicit internal and published geometric premises.
Neither script is an end-to-end formal proof of the geometry.
See \path{docs/right_angled_five_dimensional_constraints_reproduce.md}
for immutable inputs, commands and the internal/external lemma map.

\begin{thebibliography}{9}
\bibitem{Kolpakov} A. Kolpakov, \emph{On the optimality of the ideal
right-angled 24-cell}, Algebraic \& Geometric Topology 12 (2012),
1941--1960. \url{https://doi.org/10.2140/agt.2012.12.1941}.
\bibitem{Alexandrov} S. Alexandrov, \emph{On ideal vertices of right-angled
hyperbolic polyhedra}, 2023 preprint.
\url{https://arxiv.org/abs/2303.09533}.
\bibitem{Egorov} A. Egorov, \emph{On lower bounds for the number of ideal
and finite vertices of right-angled hyperbolic polyhedra in dimensions
from 5 to 12}, 2026 preprint. \url{https://arxiv.org/abs/2603.29925}.
\end{thebibliography}
\end{document}
